\documentclass[12pt]{article}
\usepackage[utf8]{inputenc}
\usepackage[T1]{fontenc}
\usepackage{hyperref}
\title{Yindjibarndi Language Resources for K-12 Curriculum Development}
\author{Compiled by Claude Sonnet 5.5 (claude-sonnet-5-5) for Andr\'as Kornai}
\date{}
\begin{document}
\maketitle
\section{Introduction}
Yindjibarndi (ISO 639-3 yij) is a Pama-Nyungan language of the Pilbara region of Western Australia, centred on Roebourne. It is not known to be taught in schools, but a local organisation is building a dictionary and grammar that it hopes will support teaching at Roebourne School.
\section{Primary Educational Resources}
No finished school textbooks were found. Juluwarlu's planned interactive multimedia dictionary and grammar is intended to support teaching at Roebourne School, and no completion date was found: \url{https://juluwarlu.com.au/our-work/language-preservation}
\section{Linguistic and Reference Materials}
The Wangka Maya Pilbara Aboriginal Language Centre page on Yindjibarndi describes the \emph{Yindjibarndi Dictionary} (2003, also viewable online as an e-book), Wordick's 1982 grammar, and recordings made between 1958 and 2001: \url{https://www.wangkamaya.org.au/pilbara-languages/yindjibarndi}

Glottolog bibliography entry (Smythe and Thieberger 1999): \url{https://glottolog.org/resource/reference/id/31027}

Endangered Languages Archive 1990 word lists (from Wangka Maya): \url{https://endangeredlanguages.com/node/94885}
\section{Religious and Community Resources}
Juluwarlu has produced many Yindjibarndi publications, and recent ones are bilingual in Yindjibarndi and English: \url{https://juluwarlu.com.au/our-work/language-preservation}

Yindjibarndi Ngurra Aboriginal Corporation RNTBC holds native title for the Yindjibarndi people. The search did not link it to any language materials: \url{https://nativetitle.org.au/node/213}
\section{Digital and Audio Resources}
Juluwarlu's planned multimedia resource would include audio, video and text examples. Archived recordings exist, per the Wangka Maya page above. No directly accessible audio collection was found.
\section{Curriculum Development Context}
This is the best-resourced of the four languages in this batch, with a dictionary, a grammar, community publications and an active local language-documentation body. The sources found do not say how the language is used in schools. Contact Juluwarlu and Wangka Maya before planning.
\end{document}